\section{Preliminaries}
\label{sec:prelims}

In the fair division problem, there is a finite set $M$ of items
that must be distributed among a finite set $N$ of agents fairly.
Formally, we are given as input a \emph{fair division instance} $\Ical \defeq (N, M, V, w)$.
Here $w \defeq (w_i)_{i \in N}$ is a collection of positive numbers that sum to 1,
and $V \defeq (v_i)_{i \in N}$ is a collection of functions,
where $v_i: 2^M \to \mathbb{R}$ and $v_i(\emptyset) = 0$ for each $i \in N$.
$v_i$ is called agent $i$'s \emph{valuation function},
and $w_i$ is called agent $i$'s \emph{entitlement}.

For any non-negative integer $k$, define $[k] \defeq \{1, 2, \ldots, k\}$.
We generally assume \wLoG{} that $N = [n]$ and $M = [m]$.
For an agent $i$ and item $j$, we often write $v_i(j)$ instead of $v_i(\{j\})$ for notational convenience.

Our task is to find a fair allocation. An \emph{allocation} $A \defeq (A_i)_{i \in N}$ is
a collection of pairwise-disjoint subsets of $M$ such that $\bigcup_{i=1}^n A_i = M$.
(If we relax the condition $\bigcup_{i=1}^n A_i = M$, we get \emph{partial allocations}.)
The set $A_i$ is called agent $i$'s \emph{bundle} in $A$.

For any function $u: 2^M \to \mathbb{R}$ and sets $S, T \subseteq M$, the \emph{marginal value}
of $S$ over $T$ is defined as $u(S \mid T) \defeq u(S \cup T) - u(T)$.

\subsection{Fairness Notions}
\label{sec:prelims:fairness-notions}

A \emph{fairness notion} $F$ is a function that takes as input a fair division instance $\Ical$,
a (partial) allocation $A$, and an agent $i$, and outputs either true or false.
When $F(\Ical, A, i)$ is true, we say that allocation $A$ is $F$\emph{-fair} to agent $i$,
or that agent $i$ is $F$\emph{-satisfied} by allocation $A$.
Allocation $A$ is said to be $F$-fair if it is $F$-fair to every agent.

A notion $F$ of fairness is said to be \emph{feasible} if for every fair division instance,
there exists an $F$-fair allocation.
We say that a notion $F_1$ of fairness \emph{implies} another notion $F_2$ of fairness if
every $F_1$-fair allocation is also an $F_2$-fair allocation.

\subsection{Efficiency}
\label{sec:prelims:efficiency}

In a fair division instance $([n], [m], (v_i)_{i=1}^n, w)$,
an allocation $A$ \emph{Pareto-dominates} an allocation $B$ if
$v_i(A_i) \ge v_i(B_i)$ for each agent $i \in [n]$,
and $v_i(A_i) > v_i(B_i)$ for some agent $i \in [n]$.
An allocation is \emph{Pareto-optimal} (PO) if it is not Pareto-dominated by any other allocation.
